\documentclass[a4paper,zihao=false]{ctexart}
\usepackage{geometry}
\geometry{left=2.5cm,right=2.5cm,top=2.5cm,bottom=2.5cm}
\usepackage{tabularx}
\usepackage{listings}

\setlength{\parindent}{2em} 
\setlength{\parskip}{0em} 

\lstset{
    basicstyle=\ttfamily,
    breaklines=true,
    frame=none,
    tabsize=4
}

\begin{document}
\fontsize{14pt}{17.5pt}\selectfont

\noindent\textbf{附件}

\noindent\textbf{1、工作日志要求（基本模板）}
\begin{center}
    \textbf{\fontsize{20pt}{25pt}\selectfont 工作日志}
\end{center}

\begin{table}[h]
\fontsize{14pt}{17.5pt}\selectfont
\renewcommand{\arraystretch}{1.5}
\begin{tabularx}{\textwidth}{|l|X|}
\hline
项目名称（内部名称） & \\ \hline
项目节点/阶段情况 & \\ \hline
研究团队成员 & \\ \hline
填报时间 & \\ \hline
填报人 & \\ \hline
\end{tabularx}
\end{table}

\noindent\textbf{一、工作目标} （当日研究工作目标）

\vspace{2em}
\noindent\textbf{二、工作内容} （当日工作内容及完成情况）

\vspace{2em}
\noindent\textbf{三、工作要点} （如果有对后续或者其它研究工作有价值的技术点、工作方法、思路等，可以作为工作要点进行记录）

\vspace{2em}
\noindent\textbf{四、待解决的问题} （当日工作遗留的问题或者项目进度确定次日需要解决的问题，可以作为次日的工作目标）





\clearpage
\noindent\textbf{2、周报要求（基本模板）}
\begin{center}
    \textbf{\fontsize{20pt}{25pt}\selectfont 研究生工作周报 }\\
    \vspace{0.5em}
    姓名：\\
    \vspace{0.5em}
    2026.2.2-2026.2.8
\end{center}

\noindent\textbf{一、本周工作目标}（本周工作目标）

\vspace{2em}
\noindent\textbf{二、本周工作情况}

\noindent\textbf{2.1、工作完成总体情况}

\vspace{1em}
\noindent\textbf{2.2、研究工作技术路线及成效}（面向具体问题，采用的技术方法、实验方法，实验效果及分析评判等

\vspace{2em}
\noindent\textbf{三、本周工作要点}（如果有对后续或者其它研究工作有价值的技术点、工作方法、思路等，可以作为工作要点进行记录）

\vspace{2em}
\noindent\textbf{四、本周工作问题总结}

\noindent\textbf{4.1、问题总结}（分别总结本周中已解决工作问题、待解决工作问题、新出现工作问题等）

\vspace{1em}
\noindent\textbf{4.2、待解决或新出现问题分析}

\vspace{1em}
\noindent\textbf{4.3、待解决或新出现问题工作思路及下一周工作计划}





\clearpage
\noindent\textbf{3、专题研究报告（基本模板）}
\begin{center}
    \textbf{\fontsize{20pt}{25pt}\selectfont 专题研究报告 }\\
    \vspace{0.5em}
    姓名：\\
\end{center}


\noindent\textbf{摘要：}

\vspace{2em}
\noindent\textbf{一、引言} （含研究背景、问题定义、技术概述、提出方法概述等）

\vspace{2em}
\noindent\textbf{二、问题建模及分析}

\vspace{2em}
\noindent\textbf{三、算法推导及分析}

\vspace{2em}
\noindent\textbf{四、仿真/验证及分析}

\vspace{2em}
\noindent\textbf{五、总结}


\vspace{2em}
\noindent\textbf{参考文献}

论文参考文献标准格式


\end{document}